\section{Domaine positif de la coordinatisation elliptique}
\begin{proposition}[Appartenance des représentations au système de coordonnées elliptique]
\label{iv:ellipse:pertenencia}
La branche positive de \eqref{eq:ii-par-angular}, avec les
intervalles \eqref{md:accion:intervalos}, vérifie
\(A>0\), \(0<C^*<A\) et \(\hbar_s>0\).
Le renversement de la coordonnée impaire conserve \(|C^*/A|<1\).
\end{proposition}
\begin{proof}
La proposition \ref{md:accion:positividad} donne
\(0<\eta_{\mathrm{ret}}<H_5\) et la positivité des deux
sections d’action. Pour une borne supérieure de \(H_5\),
\eqref{md:accion:intervalos} implique
\(\alpha<1/125\) et \(\varphi>8/5\). En conservant seulement
les termes positifs de \eqref{eq:el-h5}, on obtient
\[
 H_5<
 \frac{\sqrt5}{2}+
 \frac{9}{16\cdot125^2}+
 \frac{7}{48\cdot125^4}.
\]
Comme \(5<81/16\), on a \(\sqrt5/2<9/8\). De plus,

\[
 \frac{9}{16\cdot125^2}+
 \frac{7}{48\cdot125^4}
 =\frac{421882}{11718750000}<\frac1{1000}.
\]
Par conséquent, \(H_5<9/8+1/1000<2\). En utilisant maintenant
\(\alpha>7/1000\),
\[
 0<C^*=2(\eta_{\mathrm{ret}}+\alpha)
 <2\left(2+\frac1{125}\right)=\frac{502}{125}
 <7<1000\alpha=A.
\]
L’échange de canaux de
\ref{prop:ii-inversion-canales-angulares} remplace \(C^*\) par \(-C^*\)
et conserve \(A\) ; il conserve donc l’inégalité absolue.

\end{proof}
